
\documentclass{beamer}

%\usepackage[table]{xcolor}
\mode<presentation> {
  \usetheme{Boadilla}

%  \usetheme{Pittsburgh}
%\usefonttheme[2]{sans}
\renewcommand{\familydefault}{cmss}
%\usepackage{lmodern}
%\usepackage[T1]{fontenc}
%\usepackage{palatino}
%\usepackage{cmbright}
  \setbeamercovered{transparent}
\useinnertheme{rectangles}
}

\setbeamercolor{normal text}{fg=black}
\setbeamercolor{structure}{fg= blue}
\definecolor{trial}{cmyk}{1,0,0, 0}
\definecolor{trial2}{cmyk}{0.00,0,1, 0}
\definecolor{darkgreen}{rgb}{0,.4, 0.1}
\usepackage{array}
\beamertemplatesolidbackgroundcolor{white}  \setbeamercolor{alerted
text}{fg=red}
\usepackage{tcolorbox}
\setbeamertemplate{caption}[numbered]\newcounter{mylastframe}


\usepackage{tikz}
\usetikzlibrary{arrows}
\usepackage{colortbl}

\renewcommand{\familydefault}{cmss}
%\usepackage[all]{xy}

\usepackage{tikz}
\usepackage{lipsum}

 \newenvironment{changemargin}[3]{%
 \begin{list}{}{%
 \setlength{\topsep}{0pt}%
 \setlength{\leftmargin}{#1}%
 \setlength{\rightmargin}{#2}%
 \setlength{\topmargin}{#3}%
 \setlength{\listparindent}{\parindent}%
 \setlength{\itemindent}{\parindent}%
 \setlength{\parsep}{\parskip}%
 }%
\item[]}{\end{list}}
\usetikzlibrary{arrows}

\usecolortheme{lily}

\newtheorem{com}{Comment}
\newtheorem{lem} {Lemma}
\newtheorem{prop}{Proposition}
\newtheorem{condition}{Condition}
\newtheorem{thm}{Theorem}
\newtheorem{defn}{Definition}
\newtheorem{cor}{Corollary}
\newtheorem{obs}{Observation}
 \numberwithin{equation}{section}
 

\setbeamertemplate{navigation symbols}{}


\title[PLSC 43401]{Mathematical Foundations of Political Methodology}

\author{Robert Gulotty}
\institute[Chicago]{University of Chicago}
\vspace{0.3in}




\begin{document}

\begin{frame}
\maketitle
\end{frame}


\begin{frame}
\frametitle{Day 6: Transforms}
\begin{itemize}
\item Change of Coordinates
\item Moment Generating Functions
\end{itemize}
\end{frame}



\begin{frame}
\frametitle{Change of Coordinates}


\begin{prop}
Suppose $X$ is a random variable and $Y = g(X)$, where $g:\Re \rightarrow \Re$ that is a $monotonic$ function.  \\
Define $g^{-1}:\Re \rightarrow \Re$ such that $g^{-1}(g(X)) = X$ and is differentiable.  Then, 
\begin{eqnarray}
f_{Y}(y) & = & f_{X}(g^{-1}(y))\left|\frac{\partial g^{-1}(y)}{\partial y} \right| \text{ if $y = g(x)$ for some $x$ } \nonumber \\
& = & 0 \text{ otherwise } \nonumber 
\end{eqnarray}

\end{prop}


\end{frame}

\begin{frame}
\frametitle{Change of Coordinates}


\begin{proof}

Suppose $g(\cdot)$ is monotonically increasing (WLOG)\pause 
\begin{eqnarray}
\invisible<1>{F_{Y}(y) & = & P(Y \leq y) \nonumber \\} \pause 
\invisible<1-2>{& = & P(g(X) \leq y) \nonumber \\} \pause 
\invisible<1-3>{& = & P(X \leq g^{-1}(y) ) \nonumber \\} \pause 
\invisible<1-4>{& = & F_{X}(g^{-1}(y) ) \nonumber } \pause 
\end{eqnarray}
\invisible<1-5>{Now differentiating to get the pdf } \pause 
\begin{eqnarray}
\invisible<1-6>{\frac{\partial F_{Y}(y)}{\partial y} & = & \frac{\partial F_{X}(g^{-1}(y) )} {\partial y } \nonumber \\
 & = & f_{X}(g^{-1}(y) ) \frac{\partial g^{-1}(y)}{\partial y }   \nonumber } \pause 
\end{eqnarray}
\invisible<1-7>{Then this is a pdf because $\frac{\partial g^{-1}(Y)}{\partial y } > 0$.  } 


\end{proof}



\end{frame}


\begin{frame}
\frametitle{Change of Coordinates}

Suppose $X$ is a random variable with pdf $f_{X}(x)$.  Suppose $Y =  X^{n}$.  Find $f_{Y}(y)$. \pause  \\
\invisible<1>{Then $g^{-1} (x)  =  x^{1/n}$.  } \pause 




\begin{eqnarray}
\invisible<1-2>{f_{Y}(y) & = & f_{X}(g^{-1}(y)) \left| \frac{\partial g^{-1}(Y)}{\partial y } \right| \nonumber \\} \pause 
\invisible<1-3>{ & = & f_{X} (y^{1/n} ) \frac{y^{\frac{1}{n} - 1}}{n} \nonumber } 
\end{eqnarray}
\pause
\invisible<1-4>{\alert{We've used this to derive many of the pdfs}} \pause 
\begin{itemize}
\invisible<1-5>{\item[-] Normal distribution }
\invisible<1-5>{\item[-] Chi-Squared Distribution}
\end{itemize}




\end{frame}







\begin{frame}
\frametitle{Moment Generating Functions}


\begin{defn}
Suppose $X$ is a random variable with pdf $f$.  Define, 
\begin{eqnarray}
E[X^{n}] & = & \int_{-\infty}^{\infty} x^{n} f(x) dx   \nonumber 
\end{eqnarray}

We will call $X^{n}$ the \alert{$n^{\text{th}}$} moment of $X$ 

\end{defn}

\begin{itemize}
\item[-] By this definition $var(X) = \text{Second Moment} - \text{First Moment}^{2} $
\item[-] We are assuming that the integral converges
\end{itemize}

\end{frame}


\begin{frame}
\frametitle{Moment Generating Functions}

\begin{prop}
Suppose $X$ is a random variable with pdf $f(x)$.  Call $M(t) = E[e^{tX}]$, 
\begin{eqnarray}
M(t) & = & E[e^{tX}] \nonumber \\
 & = & \int_{-\infty}^{\infty} e^{tx} f(x) dx \nonumber 
 \end{eqnarray}

We will call $M(t)$ the moment generating function, because:
\begin{eqnarray}
\frac{\partial^{n} M (t) }{\partial^{n} t}|_{0} & = & E[X^{n}] \nonumber 
\end{eqnarray}

(Assuming that we can interchange derivative and integral)

\end{prop}



\end{frame}





\begin{frame}
\frametitle{Moment Generating Functions}

\begin{small}
\begin{proof}
Recall the Taylor Expansion of $e^{tX}$ at $0$, \pause 
\begin{eqnarray}
\invisible<1>{e^{tX} & = & 1 + tx + \frac{t^2 x^2}{2!} + \frac{t^3 x^3}{3!} + \hdots \nonumber } \pause 
\end{eqnarray}

\invisible<1-2>{Then, } \pause 
\begin{eqnarray}
\invisible<1-3>{E[e^{tX} ] & = &  1 + tE[X] + \frac{t^2}{2!} E[X^2] + \frac{t^3}{3!} E[X^3] + \hdots \nonumber } \pause 
\end{eqnarray}

\invisible<1-4>{Differentiate once:} \pause 
\begin{eqnarray}
\invisible<1-5>{\frac{\partial M(t)}{\partial t} & = & 0 + E[X]  + \frac{2t}{2!} E[X^2] + \hdots \nonumber \\
M^{'}(0) & = & 0 + E[X] + 0 + 0 \hdots \nonumber } 
\end{eqnarray}





\end{proof}

\end{small}



\end{frame}

\begin{frame}

\begin{small}
\begin{proof}

Differentiate $n$ times \pause 
\begin{eqnarray}
\invisible<1>{\frac{\partial^{n} M(t)}{\partial^{n} t } & = & 0 + 0 + 0 + \hdots + \frac{n \times n-1 \times \hdots 2 \times t^{0} E[X^{n}] }{n!}  + \frac{n!t E[X^{n+1}] }{(n+ 1)! } + \hdots \nonumber \\} \pause 
\invisible<1-2>{& = & \frac{n! E[X^{n}] }{n!}  + \frac{n!t E[X^{n+1}] }{(n+ 1)! } + \hdots  \nonumber } \pause 
\end{eqnarray}

\invisible<1-3>{Evaluated at 0, yields $M^{n}(0)  = E[X^{n}] $} \pause 

\end{proof}

\begin{itemize}
\invisible<1-4>{\item[-] If two random variables, $X$ and $Y$ have the same moment generating functions, then $F_{X}(x) = F_{Y}(y)$ for \alert{almost all} $x$.  } 
\end{itemize}


\end{small}
\end{frame}


\begin{frame}
\frametitle{The Moments of the Normal Distribution}
Suppose $Z \sim N(0,1)$.  \pause 

\begin{eqnarray}
\invisible<1>{E[e^{tX}] & = & \frac{1}{\sqrt{2\pi}} \int_{-\infty}^{\infty} e^{tx} e^{-x^2/2} dx \nonumber } \pause 
\end{eqnarray}

\invisible<1-2>{$tx - \frac{1}{2} x^2 = -\frac{1}{2}\left( ( x - t)^2 - t^2 \right)$} \pause 

\begin{eqnarray} 
\invisible<1-3>{E[e^{tX}] & = & \frac{1}{\sqrt{2\pi}}  e^{\frac{t^{2}}{2}}\int_{-\infty}^{\infty} e^{-(x- t)^2/2} dx \nonumber \\} \pause 
 \invisible<1-4>{& = & e^{\frac{t^{2}}{2}} \nonumber } 
\end{eqnarray}


\end{frame}

\begin{frame}
\frametitle{Extracting Moments of the Normal Distribution}

\pause 
\begin{eqnarray}
\invisible<1>{M^{'}(0) & = &E[X] =  e^{t^2/2} t |_{0} = 0 \nonumber \\} \pause
\invisible<1-2>{M^{''} (0 ) & = & E[X^2] =  e^{t^2/2} (t^2 + 1) |_{0} = 1 \nonumber \\} \pause 
\invisible<1-3>{M^{'''} (0) & = & E[X^3] = e^{t^2/2} t (t^2 + 3) |_{0} = 0 \nonumber \\} \pause 
\invisible<1-4>{M^{''''} (0) & = & E[X^4] = e^{t^2/2} (t^4 + 6t^2 + 3)|_{0} = 3 \nonumber } \\ \pause 
\invisible<1-5>{M^{5} (0) & = & E[X^5] = e^{t^2/2} t (t^4 + 10t^2 + 15)|_{0} = 0 \nonumber } \\ \pause 
\invisible<1-6>{M^{6} (0) & = & E[X^6] = e^{t^2/2} (t^6 + 15t^4 + 45t^2 + 15 )|_{0} = 15 \nonumber } \pause 
\end{eqnarray} 


\end{frame}

\begin{frame}

\begin{prop}
Suppose $X_{i}$ are a sequence of independent random variables.  Define 
\begin{eqnarray}
Y & = & \sum_{i=1}^{N} X_{i} \nonumber 
\end{eqnarray}

Then 
\begin{eqnarray}
M_{Y}(t) & = & \prod_{i=1}^{N} M_{X_{i}}(t) \nonumber 
\end{eqnarray}


\end{prop}

\end{frame}

\begin{frame}

\begin{proof}
\begin{eqnarray}
M_{Y}(t) & = & E[e^{tY}] \nonumber \\ \pause 
\invisible<1>{& = & E[e^{t\sum_{i=1}^{N} X_{i} } ] \nonumber \\ } \pause 
 \invisible<1-2>{& = & E[e^{tX_{1} + tX_{2} + \hdots tX_{N} } ] \nonumber \\} \pause 
  \invisible<1-3>{& = & E[e^{tX_{1} }]E[e^{tX_{2} }] \hdots E[e^{tX_{N} }] \text{ (by independence) } \nonumber \\} \pause 
   \invisible<1-4>{& = & \prod_{i=1}^{N} E[e^{tX_{i}}] \nonumber } 
\end{eqnarray}

\end{proof}

\end{frame}

\end{document}